\documentclass[12pt,a4paper]{article}
\usepackage[utf8]{inputenc}
\usepackage[margin=1in]{geometry}
\usepackage{amsmath, amssymb}
\usepackage{graphicx}
\usepackage{booktabs}
\usepackage[table]{xcolor}
\usepackage{tabularx}
\usepackage[colorlinks=true, linkcolor=blue, urlcolor=blue, citecolor=blue]{hyperref}
\usepackage{float}

\definecolor{hdrbg}{HTML}{1A237E}
\definecolor{rowalt}{HTML}{E8EAF6}
\definecolor{sigbg}{HTML}{E8F5E9}

\title{\vspace{-2cm}\color{hdrbg}\textbf{Statistical Inference (MA20266) Project}\\[1em]
\Large Topic: Does the US Federal Reserve Interest Rate Cut significantly change NIFTY 50 Stock Market Return behaviour?}
\author{\textbf{Group Members:} \\
1) Arnab Maiti (24IM10017) \\
2) Vivek Dubey (24IM10070) \\
3) Kartik Jeengar (24IM10010) \\
4) Kshitij Nayan (24IM10040)}
\date{\vspace{-1em}}

\begin{document}
\maketitle

%% ============================================================
%% PROJECT STRUCTURE
%% ============================================================
\newpage
\section{Project Structure}
\begin{table}[H]
\centering
\begin{tabularx}{\textwidth}{|p{4cm}|X|}
\hline
\rowcolor{hdrbg} \textcolor{white}{\textbf{Section}} & \textcolor{white}{\textbf{Content}} \\
\hline
1. Data Description & NIFTY 50 daily returns (2000--2026)\newline Source: Investing.com + Yahoo Finance \\
\hline
\rowcolor{rowalt}
2. Problem Statement & Do Fed rate cuts significantly change next-day NIFTY behaviour?\newline Event study: FOMC cut $\rightarrow$ next Indian trading day \\
\hline
3. Methodology & Point \& interval estimation\newline Welch's t-test (mean \& absolute returns)\newline F-test \& Jarque-Bera \\
\hline
\rowcolor{rowalt}
4. Results \& Inference & 3 of 4 hypotheses rejected\newline Statistical vs economic significance\newline Trading implications \\
\hline
\end{tabularx}
\end{table}

%% ============================================================
%% DATA DESCRIPTION
%% ============================================================
\newpage
\section{Data Description}
Dataset: Daily closing prices of the NIFTY 50 Index (\textasciicircum NSEI) from April 2000 to April 2026 --- \textbf{6457} daily observations.

Derived variables:
\begin{equation*}
    r_t = \ln\left( \frac{P_t}{P_{t-1}} \right) \times 100 \quad \text{(daily log-return)}
\end{equation*}
\begin{equation*}
    |r_t| = \left| r_t \right| \quad \text{(absolute return --- measures magnitude of price movement)}
\end{equation*}

\textbf{Event definition:} FOMC rate cut announcements happen in US evening time (IST next morning). We define the \textbf{event day} as the \textbf{next Indian trading day} after each FOMC rate cut announcement. This captures the immediate market reaction to the rate cut news.

\textbf{Source:}
\begin{itemize}
    \item Investing.com CSV download for NIFTY 50 daily prices (Apr 2000 -- May 2020)
    \item Yahoo Finance (via yfinance Python library) for Jun 2020 -- Apr 2026
    \item Federal Reserve FOMC meeting dates for rate cut event classification
\end{itemize}

\textbf{Period classification:}
\begin{table}[H]
\centering
\begin{tabularx}{\textwidth}{|l|l|X|}
\hline
\rowcolor{hdrbg} \textcolor{white}{\textbf{Category}} & \textcolor{white}{\textbf{Days}} & \textcolor{white}{\textbf{Description}} \\
\hline
Event days (E) & 31 days & Next Indian trading day after each FOMC rate cut \\
\hline
\rowcolor{rowalt}
Normal days (N) & 6426 days & All other trading days \\
\hline
\end{tabularx}
\end{table}

\textbf{Sample Data (50 observations --- all 31 event-days + 19 selected normal days):}
\begin{table}[H]
\centering
\footnotesize
\begin{tabular}{|l|c|c|c|l|c|}
\hline
\rowcolor{hdrbg}\textcolor{white}{\textbf{Date}} & \textcolor{white}{\textbf{Close}} & \textcolor{white}{\textbf{return}} & \textcolor{white}{\textbf{$|$return$|$}} & \textcolor{white}{\textbf{group}} & \textcolor{white}{\textbf{neg}} \\
\hline
04-01-2001 & 1307.65 & +1.2621 & 1.2621 & Event-day & 0 \\ \hline
01-02-2001 & 1359.15 & -0.9191 & 0.9191 & Event-day & 1 \\ \hline
16-03-2001 & 1193.55 & -1.9580 & 1.9580 & Normal-day & 1 \\ \hline
21-03-2001 & 1207.10 & +3.0405 & 3.0405 & Event-day & 0 \\ \hline
19-04-2001 & 1144.45 & +3.6528 & 3.6528 & Event-day & 0 \\ \hline
16-05-2001 & 1151.15 & +0.5095 & 0.5095 & Event-day & 0 \\ \hline
28-06-2001 & 1094.00 & -0.1918 & 0.1918 & Event-day & 1 \\ \hline
06-08-2001 & 1075.25 & +0.0605 & 0.0605 & Normal-day & 0 \\ \hline
23-08-2001 & 1071.50 & +0.2617 & 0.2617 & Event-day & 0 \\ \hline
18-09-2001 & 900.20 & +3.1541 & 3.1541 & Event-day & 0 \\ \hline
03-10-2001 & 899.65 & -1.1549 & 1.1549 & Event-day & 1 \\ \hline
02-11-2001 & 997.60 & +0.3615 & 0.3615 & Normal-day & 0 \\ \hline
07-11-2001 & 987.50 & -1.4477 & 1.4477 & Event-day & 1 \\ \hline
12-12-2001 & 1107.65 & -0.2300 & 0.2300 & Event-day & 1 \\ \hline
07-11-2002 & 960.70 & -0.1664 & 0.1664 & Event-day & 1 \\ \hline
26-06-2003 & 1116.35 & +0.8727 & 0.8727 & Event-day & 0 \\ \hline
06-09-2005 & 2428.65 & +0.2350 & 0.2350 & Normal-day & 0 \\ \hline
02-11-2006 & 3791.20 & +0.6390 & 0.6390 & Normal-day & 0 \\ \hline
12-03-2007 & 3734.60 & +0.4455 & 0.4455 & Normal-day & 0 \\ \hline
16-07-2007 & 4512.15 & +0.1686 & 0.1686 & Normal-day & 0 \\ \hline
19-09-2007 & 4732.35 & +4.0130 & 4.0130 & Event-day & 0 \\ \hline
01-11-2007 & 5866.45 & -0.5813 & 0.5813 & Event-day & 1 \\ \hline
12-12-2007 & 6159.30 & +1.0125 & 1.0125 & Event-day & 0 \\ \hline
23-01-2008 & 5203.40 & +6.0220 & 6.0220 & Event-day & 0 \\ \hline
31-01-2008 & 5137.45 & -0.5852 & 0.5852 & Event-day & 1 \\ \hline
19-03-2008 & 4573.95 & +0.8993 & 0.8993 & Event-day & 0 \\ \hline
02-05-2008 & 5228.20 & +1.1988 & 1.1988 & Event-day & 0 \\ \hline
10-10-2008 & 3279.95 & -6.8827 & 6.8827 & Event-day & 1 \\ \hline
31-10-2008 & 2885.60 & +6.7574 & 6.7574 & Event-day & 0 \\ \hline
05-12-2008 & 2714.40 & -2.6754 & 2.6754 & Normal-day & 1 \\ \hline
17-12-2008 & 2954.35 & -2.9154 & 2.9154 & Event-day & 1 \\ \hline
25-06-2010 & 5269.05 & -0.9736 & 0.9736 & Normal-day & 1 \\ \hline
02-08-2010 & 5431.65 & +1.1862 & 1.1862 & Normal-day & 0 \\ \hline
07-08-2014 & 7649.25 & -0.2976 & 0.2976 & Normal-day & 1 \\ \hline
24-08-2015 & 7809.00 & -6.0973 & 6.0973 & Normal-day & 1 \\ \hline
13-06-2016 & 8110.60 & -0.7303 & 0.7303 & Normal-day & 1 \\ \hline
06-02-2018 & 10498.25 & -1.5904 & 1.5904 & Normal-day & 1 \\ \hline
01-08-2019 & 10980.00 & -1.2490 & 1.2490 & Event-day & 1 \\ \hline
19-09-2019 & 10704.80 & -1.2611 & 1.2611 & Event-day & 1 \\ \hline
31-10-2019 & 11877.45 & +0.2812 & 0.2812 & Event-day & 0 \\ \hline
24-01-2020 & 12248.25 & +0.5559 & 0.5559 & Normal-day & 0 \\ \hline
04-03-2020 & 11251.00 & -0.4638 & 0.4638 & Event-day & 1 \\ \hline
16-03-2020 & 9197.40 & -7.9174 & 7.9174 & Event-day & 1 \\ \hline
22-06-2020 & 10311.20 & +0.6499 & 0.6499 & Normal-day & 0 \\ \hline
03-02-2022 & 17560.20 & -1.2439 & 1.2439 & Normal-day & 1 \\ \hline
05-08-2022 & 17397.50 & +0.0891 & 0.0891 & Normal-day & 0 \\ \hline
02-07-2024 & 24123.85 & -0.0750 & 0.0750 & Normal-day & 1 \\ \hline
20-09-2024 & 25790.95 & +1.4653 & 1.4653 & Event-day & 0 \\ \hline
08-11-2024 & 24148.20 & -0.2116 & 0.2116 & Event-day & 1 \\ \hline
19-12-2024 & 23951.70 & -1.0266 & 1.0266 & Event-day & 1 \\ \hline
\end{tabular}
\end{table}

%% ============================================================
%% REAL-WORLD PROBLEM
%% ============================================================
\newpage
\section{Real-World Problem}

\textbf{Research Question:} ``Does the US Federal Reserve interest rate cut significantly \textbf{change} the behaviour of NIFTY 50 returns on the next trading day?''

We investigate three dimensions of ``change in behaviour'':
\begin{enumerate}
    \item \textbf{Direction:} Do returns go up? ($\mu_E > \mu_N$)
    \item \textbf{Magnitude:} Are price movements larger? ($|r_E| > |r_N|$)
    \item \textbf{Volatility:} Is there more uncertainty? ($\sigma^2_E > \sigma^2_N$)
\end{enumerate}

\textbf{Hypothesis:} Rate cuts signal easier monetary policy globally. This creates \textbf{uncertainty} and \textbf{large price swings} in emerging markets like India, as capital flows and risk sentiment shift rapidly.

This matters for:
\begin{itemize}
    \item \textbf{Futures \& options traders:} Large event-day moves mean bigger profits/losses on leveraged positions
    \item \textbf{Risk managers:} VaR models must account for higher event-day volatility
    \item \textbf{Portfolio managers:} Position sizing around FOMC announcements
    \item \textbf{Policy economists:} Quantifying US monetary policy spillover to Indian equities
\end{itemize}

%% ============================================================
%% METHODOLOGY
%% ============================================================
\section{Methodology}

\subsection{Point Estimation}
For each group (event-day vs normal-day), we estimate:
\begin{table}[H]
\centering
\begin{tabularx}{\textwidth}{|l|X|X|}
\hline
\rowcolor{hdrbg} \textcolor{white}{\textbf{Estimator}} & \textcolor{white}{\textbf{Formula}} & \textcolor{white}{\textbf{Justification}} \\
\hline
Sample mean $\bar{x}$ & $\frac{1}{n} \sum x_i$ & Unbiased estimator of $\mu$ \\
\hline
\rowcolor{rowalt}
Sample variance $s^2$ & $\frac{1}{n-1} \sum (x_i - \bar{x})^2$ & Unbiased (Bessel's correction) \\
\hline
Mean abs return $\overline{|r|}$ & $\frac{1}{n} \sum |r_i|$ & Measures magnitude of price movement \\
\hline
\rowcolor{rowalt}
Sample proportion $\hat{p}$ & $\hat{p}_E = 16/31 = 0.516$ \newline $\hat{p}_N = 2953/6426 = 0.460$ & MLE for Bernoulli $p$ \\
\hline
\end{tabularx}
\end{table}

\subsection{Interval Estimation}
95\% Confidence Interval for the mean (t-distribution, $\sigma$ unknown):
\begin{equation*}
    \bar{x} \pm t_{\alpha/2, n-1} \cdot \frac{s}{\sqrt{n}}
\end{equation*}

95\% CI for proportion (normal approximation):
\begin{equation*}
    \hat{p} \pm z_{\alpha/2} \cdot \sqrt{\frac{\hat{p}(1-\hat{p})}{n}}
\end{equation*}

\subsection{Hypothesis Test 1 --- t-Test for Mean Returns (Direction)}
We test whether Fed rate cuts boost next-day NIFTY returns.
\begin{equation*}
    H_0: \mu_E = \mu_N \quad \text{vs} \quad H_1: \mu_E > \mu_N \quad \text{(one-tailed)}
\end{equation*}
Welch's t-statistic (unequal variances):
\begin{equation*}
    t = \frac{\bar{x}_E - \bar{x}_N}{\sqrt{\frac{s_E^2}{n_E} + \frac{s_N^2}{n_N}}}
\end{equation*}

\subsection{Hypothesis Test 2 --- t-Test for Absolute Returns (Magnitude)}
We test whether Fed rate cuts cause \textbf{bigger price movements} (regardless of direction).
\begin{equation*}
    H_0: \mu_{|E|} = \mu_{|N|} \quad \text{vs} \quad H_1: \mu_{|E|} > \mu_{|N|} \quad \text{(one-tailed)}
\end{equation*}
Same Welch's t-statistic applied to $|r_t|$ instead of $r_t$.

\subsection{Hypothesis Test 3 --- F-Test for Variance (Volatility)}
\begin{equation*}
    H_0: \sigma^2_E = \sigma^2_N \quad \text{vs} \quad H_1: \sigma^2_E \neq \sigma^2_N \quad \text{(two-tailed)}
\end{equation*}
\begin{equation*}
    F = \frac{s_E^2}{s_N^2}
\end{equation*}

\subsection{Hypothesis Test 4 --- Jarque-Bera (Normality)}
\begin{equation*}
    H_0: S = 0, K = 3 \quad \text{(normal)} \quad \text{vs} \quad H_1: \text{non-normal}
\end{equation*}
\begin{equation*}
    JB = \frac{n}{6} \left( S^2 + \frac{(K-3)^2}{4} \right)
\end{equation*}

%% ============================================================
%% RESULTS
%% ============================================================
\newpage
\section{Results \& Inference}

\textbf{Descriptive statistics --- daily log-returns (\%):}
\begin{table}[H]
\centering
\begin{tabularx}{\textwidth}{|l|c|X|}
\hline
\rowcolor{hdrbg} \textcolor{white}{\textbf{Statistic}} & \textcolor{white}{\textbf{Event-day (E)}} & \textcolor{white}{\textbf{Normal-day (N)}} \\
\hline
Sample mean $\bar{x}$ & +0.2322\% & +0.0410\% \\
\hline
\rowcolor{rowalt}
Sample std dev $s$ & 2.9730\% & 1.3409\% \\
\hline
Mean absolute return $\overline{|r|}$ & \textbf{1.9873\%} & 0.9148\% \\
\hline
\rowcolor{rowalt}
Sample size $n$ & 31 & 6426 \\
\hline
95\% CI for mean & [$-$0.8583\%, +1.3227\%] & [+0.0082\%, +0.0738\%] \\
\hline
\rowcolor{rowalt}
95\% CI for $\overline{|r|}$ & [1.1825\%, 2.7921\%] & [0.8908\%, 0.9388\%] \\
\hline
Proportion neg days $\hat{p}$ & 0.516 & 0.460 \\
\hline
\rowcolor{rowalt}
95\% CI for proportion & [0.340, 0.692] & [0.447, 0.472] \\
\hline
\end{tabularx}
\end{table}

\vspace{1em}
\textbf{Test 1 --- Welch's t-test for mean returns (direction):}
\begin{table}[H]
\centering
\begin{tabularx}{\textwidth}{|l|l|X|}
\hline
\rowcolor{hdrbg} \textcolor{white}{\textbf{Item}} & \textcolor{white}{\textbf{Value}} & \textcolor{white}{\textbf{Decision}} \\
\hline
$H_0$ & $\mu_E = \mu_N$ & \\
\hline
\rowcolor{rowalt}
$H_1$ & $\mu_E > \mu_N$ (rate cuts boost returns) & \\
\hline
t-statistic & +0.3580 & \\
\hline
\rowcolor{rowalt}
Welch df & $\sim$30 & \\
\hline
p-value (one-tailed) & 0.3614 & \textbf{Fail to Reject $H_0$} \\
\hline
\rowcolor{rowalt}
Critical value $t_{0.05}$ & +1.697 & $t < t_{crit} \Rightarrow$ fail to reject \\
\hline
\end{tabularx}
\end{table}

\textbf{Test 2 --- Welch's t-test for absolute returns (magnitude):}
\begin{table}[H]
\centering
\begin{tabularx}{\textwidth}{|l|l|X|}
\hline
\rowcolor{hdrbg} \textcolor{white}{\textbf{Item}} & \textcolor{white}{\textbf{Value}} & \textcolor{white}{\textbf{Decision}} \\
\hline
$H_0$ & $\mu_{|E|} = \mu_{|N|}$ (same magnitude) & \\
\hline
\rowcolor{rowalt}
$H_1$ & $\mu_{|E|} > \mu_{|N|}$ (bigger moves on event days) & \\
\hline
t-statistic & +2.7205 & \\
\hline
\rowcolor{rowalt}
Welch df & $\sim$30 & \\
\hline
\rowcolor{sigbg}
p-value (one-tailed) & \textbf{0.0054} & \textbf{Reject $H_0$} \\
\hline
Critical value $t_{0.05}$ & +1.697 & $t > t_{crit} \Rightarrow$ reject \\
\hline
\end{tabularx}
\end{table}

\newpage
\textbf{Test 3 --- F-test (variance / volatility):}
\begin{table}[H]
\centering
\begin{tabularx}{\textwidth}{|l|l|X|}
\hline
\rowcolor{hdrbg} \textcolor{white}{\textbf{Item}} & \textcolor{white}{\textbf{Value}} & \textcolor{white}{\textbf{Decision}} \\
\hline
$H_0$ & $\sigma^2_E = \sigma^2_N$ & \\
\hline
\rowcolor{rowalt}
$H_1$ & $\sigma^2_E \neq \sigma^2_N$ (different volatility) & \\
\hline
F-statistic & 4.9158 & \\
\hline
\rowcolor{rowalt}
Degrees of freedom & $df_1=30, df_2=6425$ & \\
\hline
\rowcolor{sigbg}
p-value (two-tailed) & \textbf{0.0000} & \textbf{Reject $H_0$} \\
\hline
\end{tabularx}
\end{table}

Volatility is \textbf{4.9$\times$ higher} on event days ($s_E = 2.97\%$ vs $s_N = 1.34\%$).

\vspace{1em}
\textbf{Test 4 --- Jarque-Bera (normality):}
\begin{table}[H]
\centering
\begin{tabularx}{\textwidth}{|l|l|X|}
\hline
\rowcolor{hdrbg} \textcolor{white}{\textbf{Item}} & \textcolor{white}{\textbf{Value}} & \textcolor{white}{\textbf{Decision}} \\
\hline
$H_0$ & Returns $\sim$ Normal ($S=0, K=3$) & \\
\hline
\rowcolor{rowalt}
$H_1$ & Returns are non-normal & \\
\hline
JB statistic & 35889.47 & \\
\hline
\rowcolor{sigbg}
p-value & 0.000000 & \textbf{Reject $H_0$} \\
\hline
Skewness ($S$) & $-$0.4977 & Left-skewed \\
\hline
\rowcolor{rowalt}
Excess Kurtosis ($K-3$) & 11.5068 & Fat tails \\
\hline
\end{tabularx}
\end{table}

\subsection{Summary of Hypothesis Tests}
\begin{table}[H]
\centering
\begin{tabularx}{\textwidth}{|l|X|c|c|}
\hline
\rowcolor{hdrbg} \textcolor{white}{\textbf{Test}} & \textcolor{white}{\textbf{Question}} & \textcolor{white}{\textbf{p-value}} & \textcolor{white}{\textbf{Decision}} \\
\hline
1. t-test (mean) & Do returns increase? & 0.3614 & Fail to Reject $H_0$ \\
\hline
\rowcolor{sigbg}
2. t-test ($|r|$) & Are moves bigger? & \textbf{0.0054} & \textbf{Reject $H_0$} \\
\hline
\rowcolor{sigbg}
3. F-test ($\sigma^2$) & Is volatility higher? & \textbf{0.0000} & \textbf{Reject $H_0$} \\
\hline
\rowcolor{sigbg}
4. Jarque-Bera & Are returns non-normal? & \textbf{0.000000} & \textbf{Reject $H_0$} \\
\hline
\end{tabularx}
\end{table}

\textbf{Result: 3 out of 4 hypotheses are rejected at $\alpha = 0.05$.} Fed rate cuts \textbf{do} significantly change NIFTY 50 behaviour.

%% ============================================================
%% STATISTICAL vs ECONOMIC SIGNIFICANCE
%% ============================================================
\section{Statistical Significance vs Economic Significance}

While the \textbf{direction} of returns is not statistically significant (Test 1), the \textbf{magnitude} and \textbf{volatility} effects are highly significant. This has major trading implications:

\begin{table}[H]
\centering
\begin{tabularx}{\textwidth}{|l|X|X|}
\hline
\rowcolor{hdrbg} \textcolor{white}{\textbf{Aspect}} & \textcolor{white}{\textbf{Statistical View}} & \textcolor{white}{\textbf{Trading Implication}} \\
\hline
Direction (mean) & Not significant ($p=0.36$) & Directionally positive ($+0.19\%$) but cannot guarantee direction \\
\hline
\rowcolor{sigbg}
Magnitude ($|r|$) & \textbf{Significant} ($p=0.0054$) & Event-day moves are $1.99\%$ vs $0.91\%$ --- \textbf{2.2$\times$ bigger} \\
\hline
\rowcolor{sigbg}
Volatility ($\sigma^2$) & \textbf{Significant} ($p \approx 0$) & Event-day $\sigma = 2.97\%$ vs $1.34\%$ --- \textbf{4.9$\times$ higher variance} \\
\hline
For Futures & $|r|$ spread $= 1.07\%$ & $\approx$ 268 NIFTY points $=$ \textbf{Rs 6,703} per lot \\
\hline
\rowcolor{rowalt}
For Options & High IV on event days & Straddle/strangle strategies profit from large moves regardless of direction \\
\hline
\end{tabularx}
\end{table}

%% ============================================================
%% CONCLUSION
%% ============================================================
\section{Overall Conclusion}

\textbf{Estimation Conclusion:} The data provides \textbf{strong evidence} that US Federal Reserve rate cuts \textbf{significantly change} NIFTY 50 return behaviour:

\begin{enumerate}
    \item \textbf{Magnitude effect (Test 2):} Price movements are \textbf{2.2$\times$ larger} on event days ($\overline{|r|}_E = 1.99\%$ vs $\overline{|r|}_N = 0.91\%$, $p = 0.0054$). This is statistically significant at $\alpha = 0.01$.
    \item \textbf{Volatility effect (Test 3):} Variance is \textbf{4.9$\times$ higher} on event days ($F = 4.92$, $p \approx 0$). Strongly significant.
    \item \textbf{Non-normality (Test 4):} Returns exhibit significant left-skew ($S = -0.50$) and fat tails ($K = 11.51$), meaning Gaussian models underestimate tail risk.
    \item \textbf{Direction (Test 1):} The mean return is directionally positive ($+0.23\%$ vs $+0.04\%$) but not statistically significant ($p = 0.36$) due to small sample size ($n = 31$) and high event-day volatility.
\end{enumerate}

\textbf{Key Insight:} Fed rate cuts \textbf{do} significantly change NIFTY 50 behaviour --- not primarily through directional returns, but through \textbf{dramatically larger price swings}. Event days see $2.2\times$ bigger moves and $4.9\times$ higher variance. This is critical for risk management, derivatives pricing, and portfolio construction around FOMC events.

\end{document}
